\section{Os objetivos de 1995}
\label{sec:objetivos-1995}

O problema da época era distribuir código por redes de máquinas
diferentes. O white paper \citetitle{paper:java-language-environment}
(maio de 1996, de Gosling e McGilton, lido hoje via Wayback) afirma que a
tecnologia deve permitir \enquote{secure, high performance, and highly
robust applications on multiple platforms in heterogeneous, distributed
networks}. O mesmo texto resume o sistema como simples, familiar,
orientado a objetos, \emph{multithreaded} e interpretado, para máxima
portabilidade e capacidades dinâmicas.

\paragraph{A decisão.}
As metas de segurança, robustez e neutralidade de arquitetura (§1.2.2 e
§1.2.3 do white paper) levaram a quatro escolhas que o resto do capítulo
retoma. Primeiro, o compilador gera \termo{bytecode}{Conjunto de
instruções da JVM, independente de processador. O \texttt{javac} o
grava em arquivos \texttt{.class}.}: um formato intermediário neutro,
pensado para transportar código a várias plataformas (§1.2.3). Segundo,
o bytecode foi desenhado para ser fácil de interpretar em qualquer
máquina ou de traduzir dinamicamente para código nativo quando o
desempenho pedir (§4.1.1). Terceiro, o código que chega pela rede não é
confiável: o sistema de execução o submete à verificação de bytecode
(§6.3), tratada na seção~\ref{sec:verificacao}. Quarto, o modelo de
memória não tem tipos ponteiro definidos pelo programador, nem aritmética
de ponteiros, e tem coleta automática de lixo (§1.2.2); o white paper
também diz que programadores Java não conseguem forjar ponteiros para a
memória (§6.1). Para a coleta, o texto descreve um \termo{garbage
collector}{Coletor de lixo: componente da JVM que libera os objetos que
o programa não alcança mais.} em segundo plano, numa thread de baixa
prioridade (§2.1.7). Já as threads são parte da linguagem, preemptivas e,
conforme a plataforma, divididas por fatias de tempo (§7.1), assunto da
seção~\ref{sec:monitores}.

A neutralidade de arquitetura tem um custo declarado: a JVM é descrita
como máquina abstrata baseada principalmente na interface POSIX
(§1.2.3), o que pressupõe suporte a multithreading na plataforma de
destino.

\paragraph{O que mudou depois.}
O formato do bytecode continuou o mesmo em espírito, mas a forma de
verificá-lo mudou (seção~\ref{sec:verificacao}) e o carregamento de
classes ficou mais elaborado (seção~\ref{sec:class-loading}). O white
paper chama o carregador de bytecode de próxima linha de defesa depois
do verificador (§6.2).

\paragraph{O que isso muda no seu código.}
Vale para qualquer JDK, inclusive o 25: você compila uma vez e o mesmo
\texttt{.class} roda onde houver JVM. Você não faz aritmética de
ponteiros nem libera objetos manualmente; em troca, o tempo de coleta e
o comportamento de threads entram no seu raciocínio sobre desempenho,
temas das seções seguintes.

\paragraph{Onde aprofundar.}
O formato do bytecode é lido em detalhe no capítulo~\ref{cap:lendo-bytecode},
e o carregamento, a ligação e a inicialização estão na
seção~\ref{sec:class-loading}.
